% Part: normal-modal-logic
% Chapter: completeness
% Section: modalities-ccs

\documentclass[../../../include/open-logic-section]{subfiles}

\begin{document}

\olfileid[tr]{nml}{com}{mod}

\olsection{Kiplikler ve Tam Tutarlı Kümeler}

\begin{explain}
Dünyalar kümesi bir normal modal mantık $\Sigma$ içindeki tam
$\Sigma$-tutarlı $\Delta$ kümelerinden oluşan bir $\mModel{M^\Sigma}$
modeli kurduğumuzda, bu ``dünyalar'' arasında bir $R^\Sigma$
erişilebilirlik bağıntısı da tanımlamalıyız. Erişilebilirlik bağıntısı (ve
$V^\Sigma$ ataması), $\mSat{M^\Sigma}{!A}[\Delta]$ ancak ve ancak
$!A\in\Delta$ olacak biçimde tanımlansın istiyoruz. Bunu nasıl yapmalıyız?

Erişilebilirlik bağıntısı tanımlandığında, bir dünyadaki doğruluğun tanımı
\iftag{prvBox}{$R^\Sigma\Delta\Delta'$ olan bütün $\Delta'$ için
  $\mSat{M^\Sigma}{!A}[\Delta']$ ancak ve ancak
  $\mSat{M^\Sigma}{\Box!A}[\Delta]$ olduğunu}{$R^\Sigma\Delta\Delta'$
  olan en az bir $\Delta'$ için $\mSat{M^\Sigma}{!A}[\Delta']$ ancak ve
  ancak $\mSat{M^\Sigma}{\Diamond!A}[\Delta]$ olduğunu} güvence altına
alır. $\mSat{M^\Sigma}{!A}[\Delta]$ ancak ve ancak $!A\in\Delta$
olduğunun ispatı, bunun özellikle bir kiplik işleciyle başlayan !!{formula}s
için de doğru olmasını gerektirir; yani
\iftag{prvBox}{$\mSat{M^\Sigma}{\Box!A}[\Delta]$ ancak ve ancak
  $\Box!A\in\Delta$}{$\mSat{M^\Sigma}{\Diamond!A}[\Delta]$ ancak ve ancak
  $\Diamond!A\in\Delta$}. Bu gerekliliği
\iftag{prvBox}{$\Box !A$}{$\Diamond !A$} için bir dünyadaki doğruluk
tanımıyla birleştirirsek şunu elde ederiz:
\[
\iftag{prvBox}{%
  \Box!A \in \Delta \text{ ancak ve ancak } !A \in \Delta' \text{, }
    R^\Sigma\Delta\Delta' \text{ olan bütün } \Delta' \text{ için}}{%
  \Diamond!A \in \Delta \text{ ancak ve ancak } !A \in \Delta' \text{, }
    R^\Sigma\Delta\Delta' \text{ olan en az bir } \Delta' \text{ için}}
\]
\iftag{prvBox}{Soldan sağa yönü ele alın: $\Box!A\in\Delta$ ise her $!A$
  ve $R^\Sigma\Delta\Delta'$ olan her $\Delta'$ için $!A\in\Delta'$
  olduğunu söyler. $R^\Sigma\Delta\Delta'$ ancak ve ancak
  $\Delta$ içindeki bütün $\Box!A$ formülleri için $!A\in\Delta'$ olacak
  biçimde tanımlanırsa bu geçerli olur. ``Ancak ve ancak''ın sağındaki
  koşulu daha kısa olarak
  $\Setabs{!A}{\Box!A\in\Delta}\subseteq\Delta'$ yazabiliriz.}{Sağdan
  sola yönü ele alın: $!A\in\Delta'$ ise, her $!A$ ve
  $R^\Sigma\Delta\Delta'$ olan her $\Delta'$ için
  $\Diamond!A\in\Delta$ olduğunu söyler. $\Delta'$ içindeki bütün $!A$
  formülleri için $\Diamond!A\in\Delta$ olduğunda
  $R^\Sigma\Delta\Delta'$ geçerli olacak biçimde tanımlarsak bu sonuç
  geçerli olur. ``Ancak ve ancak''ın sağındaki koşulu daha kısa olarak
  $\Setabs{\Diamond!A}{!A\in\Delta'}\subseteq\Delta$ yazabiliriz.}

Soru şudur: $R^\Sigma$'nın bu tanımı gerçekten
\iftag{prvBox}{$\Box!A\in\Delta$ ancak ve ancak
  $\mSat{M^\Sigma}{\Box!A}[\Delta]$ olmasını güvence altına alır mı?
  \iftag{prvDiamond}{Ayrıca şunu güvence altına alır mı:}{}{}}{}
\iftag{prvDiamond}{$\Diamond!A\in\Delta$ ancak ve ancak
  $\mSat{M^\Sigma}{\Diamond!A}[\Delta]$?}{} Sonraki birkaç sonuç bunu
kuracaktır.
\end{explain}

\begin{defn}
  $\Gamma$ bir !!{formula}s kümesiyse
  \begin{align*}
    \Box\Gamma & = \Setabs{\Box !B}{!B \in \Gamma}\\
    \Diamond\Gamma & = \Setabs{\Diamond !B}{!B \in \Gamma}\\
    \intertext{ve}
    \Box^{-1}\Gamma & = \Setabs{!B}{\Box !B \in \Gamma}\\
    \Diamond^{-1}\Gamma & = \Setabs{!B}{\Diamond !B \in \Gamma}\\
  \end{align*}
  olsun.
\end{defn}

Başka bir deyişle $\Box\Gamma$, $\Gamma$ içindeki her !!{formula}{}ün önüne
$\Box$ konmuş hâlidir; $\Box^{-1}\Gamma$ ise $\Gamma$ içindeki bütün
$\Box$'lı !!{formula}s{}in baştaki $\Box$'ları kaldırılmış hâlidir. Bu tanım
tek başına çok önemli değildir; fakat gösterimi epeyce sadeleştirecektir.

$\Box\Box^{-1}\Gamma\subseteq\Gamma$ olduğuna dikkat edin:
\[
\Box\Box^{-1}\Gamma = \Setabs{\Box!B}{\Box!B \in \Gamma}
\]
yani bu, $\Gamma$ içinde $\Box$ ile başlayan bütün !!{formula}s{}in kümesidir.

\begin{lem}\ollabel{lem:box1}
  $\Gamma\Proves[\Sigma]!A$ ise
  $\Box\Gamma\Proves[\Sigma]\Box!A$.
\end{lem}

\begin{proof}
  $\Gamma\Proves[\Sigma]!A$ ise $!B_1,\dots,!B_k\in\Gamma$ bulunur ve
  $\Sigma\Proves !B_1\lif(!B_2\lif\cdots(!B_k\lif !A)\cdots)$ olur.
  $\Sigma$ normal olduğundan \RK{} kuralıyla
  $\Sigma\Proves\Box!B_1\lif(\Box!B_2\lif\cdots(\Box!B_k\lif
  \Box!A)\cdots)$; açıkça $\Box!B_1,\dots,\Box!B_k\in\Box\Gamma$.
  Dolayısıyla tanım gereği $\Box\Gamma\Proves[\Sigma]\Box!A$.
\end{proof}

\begin{lem}\ollabel{lem:box2}
  $\Box^{-1}\Gamma\Proves[\Sigma]!A$ ise
  $\Gamma\Proves[\Sigma]\Box!A$.
\end{lem}

\begin{proof}
  $\Box^{-1}\Gamma\Proves[\Sigma]!A$ olsun. \olref{lem:box1} gereğince
  $\Box\Box^{-1}\Gamma\Proves[\Sigma]\Box!A$. Fakat
  $\Box\Box^{-1}\Gamma\subseteq\Gamma$ olduğundan tekdüzelik gereğince
  $\Gamma\Proves[\Sigma]\Box!A$.
\end{proof}

\iftag{prvBox}{%
% Box is primitive

\begin{prop}\ollabel{prop:box}
  $\Gamma$ tam $\Sigma$-tutarlıysa, $\Box!A\in\Gamma$ olması,
  $\Box^{-1}\Gamma\subseteq\Delta$ olan her tam $\Sigma$-tutarlı
  $\Delta$ için $!A\in\Delta$ olmasına denktir.
\end{prop}

\begin{proof}
  $\Gamma$ tam $\Sigma$-tutarlı olsun. ``Yalnızca eğer'' yönü kolaydır:
  $\Box!A\in\Gamma$ ve $\Box^{-1}\Gamma\subseteq\Delta$ olsun.
  $\Box!A\in\Gamma$ olduğundan $!A\in\Box^{-1}\Gamma\subseteq\Delta$;
  dolayısıyla $!A\in\Delta$.

  ``Eğer'' yönü için karşıt-tersi ispatlayalım. $\Box!A\notin\Gamma$ olsun.
  $\Gamma$ tam $\Sigma$-tutarlı olduğundan türetim altında kapalıdır;
  dolayısıyla $\Gamma\Proves/[\Sigma]\Box!A$. \olref{lem:box2} gereğince
  $\Box^{-1}\Gamma\Proves/[\Sigma]!A$. Bundan
  \olref[prf][con]{prop:consistencyfacts}\olref[prf][con]{prop:consistencyfacts-b}
  gereğince $\Box^{-1}\Gamma\cup\{\lnot!A\}$ $\Sigma$-tutarlıdır.
  Lindenbaum Lemmasıyla, $\Box^{-1}\Gamma\cup\{\lnot!A\}\subseteq\Delta$
  olacak bir tam $\Sigma$-tutarlı $\Delta$ vardır. Tutarlılık gereğince
  $!A\notin\Delta$.
\end{proof}
}{%
% Box is not primitive, only Diamond is

\begin{prop}\ollabel{prop:diamond}
  $\Gamma$ tam $\Sigma$-tutarlıysa, $\Diamond!A\in\Gamma$ olması,
  $\Diamond\Delta\subseteq\Gamma$ ve $!A\in\Delta$ olacak tam
  $\Sigma$-tutarlı bir $\Delta$ bulunmasına denktir.
\end{prop}

\begin{proof}
  $\Gamma$ tam $\Sigma$-tutarlı olsun. Sağdan sola yön kolaydır:
  $\Diamond\Delta\subseteq\Gamma$ ve $!A\in\Delta$ olan tam
  $\Sigma$-tutarlı bir $\Delta$ alalım. O zaman
  $\Diamond!A\in\Diamond\Delta\subseteq\Gamma$.

  Soldan sağa yön için $\Diamond!A\in\Gamma$ varsayalım. $!A\in\Delta$ ve
  $\Diamond\Delta\subseteq\Gamma$ olan tam $\Sigma$-tutarlı bir $\Delta$
  bulunduğunu göstermeliyiz. $\Diamond!A\in\Gamma$ ve $\Gamma$
  $\Sigma$-tutarlı olduğundan $\Box\lnot!A\notin\Gamma$. $\Gamma$ tam
  $\Sigma$-tutarlı ve türetim altında kapalı olduğundan
  $\Gamma\Proves/[\Sigma]\Box\lnot!A$. \olref{lem:box2} gereğince
  $\Box^{-1}\Gamma\Proves/[\Sigma]\lnot!A$. Bundan
  \olref[prf][con]{prop:consistencyfacts}\olref[prf][con]{prop:consistencyfacts-b}
  gereğince $\Box^{-1}\Gamma\cup\{!A\}$ $\Sigma$-tutarlıdır. Lindenbaum
  Lemmasıyla $\Box^{-1}\Gamma\cup\{!A\}\subseteq\Delta$ olacak bir tam
  $\Sigma$-tutarlı $\Delta$ vardır. Açıkça $!A\in\Delta$. Şimdi
  $\Diamond\Delta\subseteq\Gamma$ olduğunu gösterelim. Olmadığını
  varsayalım: bazı $!B\in\Delta$ için $\Diamond!B\notin\Gamma$. O zaman
  $\Gamma$ tam $\Sigma$-tutarlı olduğundan $\Box\lnot!B\in\Gamma$.
  Dolayısıyla $\lnot!B\in\Box^{-1}\Gamma\subseteq\Delta$ olurdu ve bu da
  $\Delta$'yı tutarsız yapardı.
\end{proof}
}

\begin{lem}\ollabel{lem:box-iff-diamond}
  $\Gamma$ ve $\Delta$ tam $\Sigma$-tutarlı olsun. O zaman
  $\Box^{-1}\Gamma\subseteq\Delta$ ancak ve ancak
  $\Diamond\Delta\subseteq\Gamma$.
\end{lem}

\begin{proof}
  ``Yalnızca eğer'' yönü: $\Box^{-1}\Gamma\subseteq\Delta$ varsayalım ve
  $\Diamond!A\in\Diamond\Delta$ (yani $!A\in\Delta$) olsun.
  $\Diamond!A\in\Gamma$ olduğunu göstermek için
  $\Box\lnot!A\notin\Gamma$ olduğunu göstermek yeterlidir; çünkü o zaman
  azamilik gereğince $\lnot\Box\lnot!A\in\Gamma$. Eğer
  $\Box\lnot!A\in\Gamma$ olsaydı varsayım gereğince
  $\lnot!A\in\Delta$ olur ve bu $!A\in\Delta$ olduğundan $\Delta$'nın
  tutarlılığına aykırı düşerdi. Dolayısıyla istenildiği gibi
  $\Box\lnot!A\notin\Gamma$.

  ``Eğer'' yönü: $\Diamond\Delta\subseteq\Gamma$ varsayalım. Karşıt-ters
  yönde, $!A\notin\Delta$ varsayımından $\Box!A\notin\Gamma$ sonucunu
  gösterelim. $!A\notin\Delta$ ise azamilik gereğince
  $\lnot!A\in\Delta$; varsayımdan $\Diamond\lnot!A\in\Gamma$. Normal
  modal mantıkta $\Diamond\lnot!A$, $\lnot\Box!A$'ya denktir; ikincisi
  $\Gamma$ içindeyse tutarlılık gereğince $\Box!A\notin\Gamma$.
\end{proof}

\iftag{prvBox}{%
  \iftag{prvDiamond}{%
% Box and Diamond are both primitive; prove prop:diamond

\begin{prop}\ollabel{prop:diamond}
  $\Gamma$ tam $\Sigma$-tutarlıysa, $\Diamond!A\in\Gamma$ olması,
  $\Diamond\Delta\subseteq\Gamma$ ve $!A\in\Delta$ olacak tam
  $\Sigma$-tutarlı bir $\Delta$ bulunmasına denktir.
\end{prop}

\begin{proof}
$\Gamma$ tam $\Sigma$-tutarlı olsun. \Dual{} ve kapalılık gereğince
$\Diamond!A\in\Gamma$ ancak ve ancak $\lnot\Box\lnot!A\in\Gamma$.
$\Gamma$ tam $\Sigma$-tutarlı olduğundan
\olref[ccs]{prop:ccs-properties}\olref[ccs]{prop:ccs-lnot} gereğince
$\lnot\Box\lnot!A\in\Gamma$ ancak ve ancak
$\Box\lnot!A\notin\Gamma$. \olref{prop:box} gereğince
$\Box\lnot!A\notin\Gamma$ olması,
$\Box^{-1}\Gamma\subseteq\Delta$ ve $\lnot!A\notin\Delta$ olacak tam
$\Sigma$-tutarlı bir $\Delta$ bulunmasına denktir. Böyle bir $\Delta$ ele
alalım. \olref{lem:box-iff-diamond} gereğince
$\Box^{-1}\Gamma\subseteq\Delta$ ancak ve ancak
$\Diamond\Delta\subseteq\Gamma$. Ayrıca
\olref[ccs]{prop:ccs-properties}\olref[ccs]{prop:ccs-lnot} gereğince
$\lnot!A\notin\Delta$ ancak ve ancak $!A\in\Delta$. Dolayısıyla
$\Diamond!A\in\Gamma$ olması, $\Diamond\Delta\subseteq\Gamma$ ve
$!A\in\Delta$ olacak tam $\Sigma$-tutarlı bir $\Delta$ bulunmasına denktir.
\end{proof}

}{}}{}

\begin{prob}
  $\Gamma$ tam $\Sigma$-tutarlıysa
  \iftag{prvBox}{$\Diamond!A\in\Gamma$ olmasının,
    $\Box^{-1}\Gamma\subseteq\Delta$ ve $!A\in\Delta$ olacak bir tam
    $\Sigma$-tutarlı $\Delta$ bulunmasına denk olduğunu}{$\Box!A\in\Gamma$
    olmasının, $\Box^{-1}\Gamma\subseteq\Delta$ olan her tam
    $\Sigma$-tutarlı $\Delta$ için $!A\in\Delta$ olmasına denk olduğunu}
  gösterin. Bunu \emph{\olref[nml][com][mod]{lem:box-iff-diamond} lemmasını
  kullanmadan} yapın.
\end{prob}

\end{document}
